% =============================================================================
\section{Reconstrução, retentor de ordem zero e a identidade de Poisson}
% =============================================================================

\subsection*{Teoria mínima}

\paragraph{Reconstrução ideal.} O filtro de reconstrução é um passa-baixas
ideal de ganho $T_s$ e corte $F_s/2$, cuja resposta ao impulso é um sinc:
\[
  h_r(t)=\operatorname{sinc}\!\Bigl(\frac{t}{T_s}\Bigr),\qquad
  x(t)=\sum_n x[n]\,\operatorname{sinc}\!\Bigl(\frac{t-nT_s}{T_s}\Bigr).
\]

\paragraph{Retentor de ordem zero (ZOH).} Na prática usa-se
$h_{zoh}(t)=1$ para $0\le t<T_s$:
\[
  H_{zoh}(j\Omega)=T_s\,e^{-j\Omega T_s/2}\,
  \operatorname{sinc}\!\Bigl(\frac{f}{F_s}\Bigr),
\]
com duas consequências: (i) um \textbf{atraso} de $T_s/2$; (ii) uma
\textbf{atenuação} que chega a $-3{,}92$ dB em $f=F_s/2$ (queda de sinc).

\paragraph{A identidade de Poisson — e o termo que quase todo mundo esquece.}
Tomando $h[n]=T_s\,h_a(nT_s)$ (a definição da invariância ao impulso do
Capítulo~\ref{sec:invimp}), vale a soma periódica
\[
  H(e^{j\omega})=\sum_{k=-\infty}^{\infty}
  H_a\!\left(j\frac{\omega-2\pi k}{T_s}\right)
  \;+\;\underbrace{\frac{T_s\,h_a(0^+)}{2}}_{\text{meia-descontinuidade}}
\]
O último termo é necessário sempre que $h_a(t)$ tem \textbf{salto em
$t=0$} (em particular quando $h_a(0^+)\neq0$ e $h_a(0^-)=0$, o caso de
\emph{todos} os filtros causais analógicos). Sem ele a soma de réplicas fica
exatamente $T_s h_a(0^+)/2$ abaixo do valor verdadeiro.

\ideia{Duas leituras práticas: (1) a resposta do filtro digital é a soma de
``cópias'' do protótipo analógico deslocadas de $kF_s$ — logo, baixando $F_s$,
as cópias invadem a banda-base; (2) o termo de meia-descontinuidade explica
por que o ganho em DC do filtro do desafio não é zero, mas $\approx T_s/2$.}

\subsection*{Exercícios resolvidos}

% ---------------------------------------------------------------- 7.1
\exer{Um ZOH reconstrói um sinal amostrado a $F_s$. Calcule a atenuação do
filtro de retenção em $f=F_s/2$ e em $f=0{,}4F_s$, e o atraso introduzido.}

\textbf{Solução.} $|H_{zoh}|=T_s|\operatorname{sinc}(f/F_s)|$:
\[
  f=\frac{F_s}{2}:\quad |\operatorname{sinc}(0{,}5)|=\frac{\operatorname{sen}(\pi/2)}{\pi/2}
  =\frac{2}{\pi}=0{,}636620\ \Rightarrow\
  -3{,}9224\ \text{dB};
\]
\[
  f=0{,}4F_s:\quad |\operatorname{sinc}(0{,}4)|=0{,}756827\ \Rightarrow\
  -2{,}4201\ \text{dB}.
\]
O atraso de fase é $\phi=-\Omega T_s/2\Rightarrow$ atraso de $T_s/2$ (meia
amostra) em \emph{todas} as frequências.

\resp{$-3{,}9224$ dB em $F_s/2$; $-2{,}4201$ dB em $0{,}4F_s$; atraso de
$T_s/2$. É por isso que conversores D/A práticos usam ``compensação de
$\operatorname{sinc}$''.}

% ---------------------------------------------------------------- 7.2
\exer{Verifique numericamente a identidade de Poisson (com e sem a correção de
meia-descontinuidade) para $h_a(t)=e^{-1000t}u(t)$ com $T_s=0{,}1$ ms, na
frequência $\omega=0$.}

\textbf{Solução.} O valor exato vem da transformada $Z$ da sequência amostrada
(polos reais):
\[
  H(z)=\frac{T_s}{1-e^{-aT_s}z^{-1}},\qquad
  H(e^{j0})=\frac{T_s}{1-e^{-aT_s}}
  =\frac{10^{-4}}{1-0{,}9048374180}=1{,}05083319\times10^{-3}.
\]
Somando as réplicas $\sum_k\frac{1}{a+j(\omega-2\pi k)/T_s}$ obtém-se, sem
correção, $1{,}00083307\times10^{-3}$; com o termo
$T_sh_a(0^+)/2=T_s/2=5\times10^{-5}$ obtém-se
$1{,}05083307\times10^{-3}$ — que é o valor exato até a 7ª casa.
Na varredura de frequência, o desvio médio das réplicas ``puras'' em relação
ao exato é $-5{,}00001266\times10^{-5}\approx-T_s/2$.

\resp{Sem a correção falta exatamente $T_s h_a(0^+)/2=5\times10^{-5}$; com ela
o erro relativo cai para $\sim10^{-7}$.}

% ---------------------------------------------------------------- 7.3
\exer{\textbf{(o filtro do Desafio 08)} Explique e comprove numericamente por
que $\sum_n h[n]\approx T_s/2$ e não zero, apesar de $H_a(j0)=0$.}

\textbf{Solução.} Para o filtro do desafio, $H_a(s)=s/(s^2+200\pi s+\cdots)$
tem $H_a(j0)=0$; portanto, em princípio, a soma dos $h[n]$ deveria ser nula
(do mesmo modo que $\int h_a(t)dt=H_a(0)=0$). Duas razões impedem o zero
exato:
\begin{enumerate}
  \item $\sum_n h[n]=T_s\sum_n h_a(nT_s)$ é uma \emph{soma de Riemann} de
        $\int h_a=0$ — aproximação, não igualdade;
  \item a identidade de Poisson exige o termo
        $T_sh_a(0^+)/2=T_s/2$, pois $h_a(0^+)=1$ (salto em $t=0$).
\end{enumerate}
O valor medido é
\[
  \left|H(e^{j0})\right|=1{,}2533725398\times10^{-5}
  \quad\text{contra}\quad \frac{T_s}{2}=1{,}25\times10^{-5},
\]
isto é, $+0{,}27\%$ — praticamente \textbf{metade da primeira amostra}
($h[0]=T_s=2{,}5\times10^{-5}$).

\resp{$\sum_n h[n]=|H(e^{j0})|=1{,}25337\times10^{-5}\approx T_s/2$: o resíduo
não é erro numérico, é o termo de meia-descontinuidade (o filtro é
\emph{causal}, logo tem salto em $t=0$).}
